\section{Method}

We present \method, a split-context prompting method with three phases that directly address the challenges identified in Section~\ref{sec:intro}.
\emph{Planning} decomposes requests into evidence-specific conditions, handling structured text and compositional logic.
\emph{Adaptation} detects heterogeneity and adversarial content before execution.
\emph{Execution} aggregates evidence using three-valued logic with correct uncertainty propagation.
Planning runs once per request and is cached; only Adaptation and Execution run per item, requiring $O(k)$ LLM calls where $k$ is the number of conditions.

\subsection{Phase 1: Planning}

Planning transforms a natural-language request into a validated, executable script.
Given the user request, the LLM first generates a structured plan identifying: (1) each condition to verify, (2) the evidence source for each condition (\texttt{ATTR}, \texttt{HOURS}, \texttt{REVIEW\_TEXT}, \texttt{REVIEW\_META}), (3) value interpretation rules, and (4) the compositional aggregation logic.

The plan is then translated into a \emph{Lightweight Workflow Template} (LWT), a simple script where each line specifies an LLM call with explicit data access paths:
\begin{verbatim}
(0)=LLM("Check {(input)}[attr][NoiseLevel]. Quiet?
        Output: 1=yes, -1=no, 0=unknown")
(1)=LLM("Check {(input)}[attr][WiFi]. Free?
        Output: 1=yes, -1=no, 0=unknown")
(2)=LLM("Analyze {(input)}[reviews][0..N][text].
        Mention good matcha? Output: 1/-1/0")
(3)=LLM("Results: Noise={(0)}, WiFi={(1)}, Matcha={(2)}.
        Logic: AND. Output: -1 if any -1, 1 if all 1, else 0")
\end{verbatim}

Both plan and script undergo iterative validation (up to 5 rounds) to ensure completeness and schema compliance.
Validated scripts are cached per request, amortizing planning cost across all candidate items.

\subsection{Phase 2: Adaptation}

Adaptation analyzes actual data content to detect characteristics requiring workflow modification.
This phase runs per item, as content varies across candidates.

The LLM examines review content for three potential issues:
\begin{itemize}
    \item \textbf{Heterogeneity}: Highly uneven review lengths where verbose reviews might dominate.
    \item \textbf{Attack patterns}: Command-like text (``ignore previous'', ``output:'', ``answer is'').
    \item \textbf{Fake indicators}: Generic, suspiciously uniform reviews covering all aspects.
\end{itemize}

Based on detected conditions, the workflow may be modified---e.g., adding a filtering step that excludes suspicious content before the main analysis, or weighting reviews by credibility signals.
Critically, adaptation happens \emph{before} execution: the main reasoning phase never sees potentially adversarial content, only the controlled detection phase does.
This architectural isolation provides robustness without degrading performance on clean data.

\subsection{Phase 3: Execution}

Execution runs the (possibly adapted) script step-by-step.
Variable references are resolved dynamically: \texttt{\{(input)\}} becomes the current item's data, and \texttt{\{(N)\}} becomes the output of step $N$.
Each evidence source is processed independently, preventing cross-contamination between metadata checks and review analysis.

The final step aggregates evidence using three-valued logic (Section~\ref{sec:three-valued}), ensuring correct compositional reasoning: for AND, any violation dominates; for OR, any satisfaction dominates; unknown values propagate only when they could affect the outcome.
This preserves minority viewpoints---a single negative review is not drowned out by neutral ones.
